% TODO: kalimat tentang CO (nol dari lima uji, estimasi tertinggi 27 ppm, belum dikalibrasi) bersifat sementara.
% Perbarui setelah pengembangan prototipe CO selesai.
\begin{onehalfspace}
Kabin kendaraan yang tertutup dapat menumpuk karbon dioksida (CO\textsubscript{2}) dari napas penumpang atau kemasukan karbon monoksida (CO) dari kendaraan lain, sedangkan pengemudi sulit menyadarinya. Dokumen C-251 merencanakan AEROS (\textit{Automatic Environmental Response and On-board Safety}), yaitu prototipe pemantau CO dan CO\textsubscript{2} yang memberi peringatan serta membuka jendela otomatis. Pada C-501, rancangan diimplementasikan dengan ESP32. Sensor CO memakai MiCS-5524 yang dibaca lewat ADS1115, sedangkan sensor CO\textsubscript{2} memakai MH-Z19C yang berprinsip NDIR. Ambang peringatan ditetapkan dari acuan paparan serta data rekaman. Ambang CO\textsubscript{2} adalah 2.500 ppm untuk waspada dan 4.000 ppm untuk bahaya. Ambang CO adalah 35 ppm untuk waspada dan 87 ppm untuk bahaya. Peringatan berupa bunyi \textit{buzzer} serta sinyal \textit{relay} yang menurunkan jendela saat bahaya. Data dikirim ke basis data awan lewat MQTT dengan TLS dan HMAC. Pengujian dilakukan lima kali di tiga mobil dengan empat penumpang dan jendela tertutup. Kadar CO\textsubscript{2} naik dari 2.500 ke 4.000 ppm dalam 7--15 menit pada empat uji.  Deteksi CO belum terbukti (nol dari lima uji) karena tidak ada sumber CO acuan. Estimasi CO tertinggi di mobil hanya 27 ppm. Nilai CO bersifat estimasi karena sensor tidak selektif dan belum dikalibrasi dengan alat referensi. Jendela masih digerakkan manual sehingga yang terbukti hanya sinyal aktuasinya. Prototipe ini berfungsi sebagai peringatan dini CO\textsubscript{2}. Kalibrasi dengan gas acuan dan pengujian aktuator pada jendela sebenarnya menjadi saran lanjutan.

\noindent \textbf{Kata Kunci:} Kualitas Udara Kabin, Karbon Monoksida, Karbon Dioksida, \textit{Internet of Things}, ESP32, MQTT, Pemantauan \textit{Real-Time}, Peringatan Dini, AEROS
\end{onehalfspace}
